\section{Methodology}
\label{sec:methodology}

The GeoSoil-Cuba workflow integrates soil-map
preparation, vegetation and land-cover information,
and Sentinel-2 satellite observations through a
sequence of spatial processing and analytical
stages. The methodology addresses soil-map
coverage and topology, integration of soil and
vegetation attributes, satellite-image processing,
and temporal analysis of the Normalized Difference
Vegetation Index (NDVI).

\subsection{Workflow Overview}
\label{subsec:workflow-overview}

The main processing stages and their sequence
are illustrated in \cref{fig:workflow}.

\begin{figure}[H]
	\centering
	\begin{tikzpicture}[
	node distance=0.65cm,
	workflow/.style={
		rectangle,
		rounded corners=3pt,
		draw=geosoilgreen,
		fill=geosoilgreen!8,
		line width=0.8pt,
		text width=10cm,
		minimum height=1.05cm,
		align=center,
		font=\small
	},
	arrow/.style={
		-{Stealth[length=3mm]},
		thick,
		draw=geosoilgreen
	}
	]
	
	\node[workflow] (soil) {
		\textbf{Stage 1: Soil Data Preparation}\\
		CRS harmonization, geometry validation,
		AOI clipping, and soil coverage
	};
	
	\node[workflow, below=of soil] (vegetation) {
		\textbf{Stage 2: Soil--Vegetation Integration}\\
		Land-cover classification and spatial analysis
	};
	
	\node[workflow, below=of vegetation] (acquisition) {
		\textbf{Stage 3: Sentinel-2 Data Acquisition}\\
		Product discovery, band acquisition,
		and quality control
	};
	
	\node[workflow, below=of acquisition] (ndvi) {
		\textbf{Stage 4: NDVI Processing}\\
		Reflectance correction, NDVI calculation,
		mosaicking, and spatial extraction
	};
	
\node[workflow, below=of ndvi] (temporal) {
	\textbf{Stage 5: Temporal NDVI Analysis}\\
	Regional statistics, QGIS-based field sampling,
	time series, and visualization
};
	
	\draw[arrow] (soil) -- (vegetation);
	\draw[arrow] (vegetation) -- (acquisition);
	\draw[arrow] (acquisition) -- (ndvi);
	\draw[arrow] (ndvi) -- (temporal);
	
\end{tikzpicture}
	\caption{
		Simplified GeoSoil-Cuba workflow showing
		the principal stages of spatial data
		preparation, integration, and NDVI analysis.
	}
	\label{fig:workflow}
\end{figure}

The workflow comprises five principal stages:
soil-data inspection and preparation; soil-map
intersection with the Area of Interest (AOI);
integration of soil and vegetation information;
acquisition and quality screening of Sentinel-2
imagery; and calculation and analysis of NDVI.

The stages were implemented through a series
of Python notebooks, with QGIS used for the
visual selection of agricultural sampling areas.
Intermediate spatial datasets, tables, and
figures were exported to support quality control,
reproducibility, and subsequent analysis.


\subsection{Soil Data Preparation and AOI Intersection}
\label{subsec:soil-preparation}

The Mayabeque soil map at a scale of 1:25,000
was inspected before spatial analysis. The soil
dataset and the AOI were evaluated for their
coordinate reference systems and geometry
characteristics.

The AOI, originally supplied in WGS 84 geographic
coordinates (EPSG:4326), was transformed to
NAD27 / Cuba Norte (EPSG:3795), the projected
coordinate reference system used for soil-map
analysis.

Soil geometries were checked for validity.
Invalid geometries were repaired using
\texttt{make\_valid()}, and resulting geometry
collections were normalized to polygonal
geometries where necessary. The original source
data were preserved, and repaired geometries
were stored separately for analysis.

Soil features intersecting the AOI were selected
and clipped to its boundary. The resulting
dataset was used to assess soil-map coverage
and calculate the areas represented by mapped
soil units.

Spatial differences between the AOI and the
mapped soil coverage were examined alongside
the water-body dataset to distinguish
reservoir-related areas from other uncovered
locations.

Areas not represented by soil polygons and not
explained by the mapped reservoir were retained
as unresolved data gaps. No soil classifications
were assigned to these areas without supporting
information.

The resulting soil dataset was also examined
for undefined classification attributes and
overlapping polygon geometries. These conditions
were documented as data-quality findings.
Positive-area overlaps were identified but
were not automatically eliminated by modifying
the original soil-unit boundaries.


\subsection{Soil--Vegetation Integration}
\label{subsec:soil-vegetation-integration}

The processed soil layer was integrated with
existing vegetation and land-cover information
to support spatial analysis of mapped soil units
in their vegetation context.

The attributes \texttt{LCNAME}, \texttt{NEWVEG},
and \texttt{TNC\_CLASS} were retained for
analysis. Missing or undefined classification
values were identified during data preparation
rather than being replaced with unsupported
categories.

The integrated vector dataset provides a common
spatial framework for examining the distribution
of vegetation classifications across mapped
soil units. It also supports the preparation
of thematic maps and the examination of spatial
coverage and classification completeness.

The resulting dataset was exported in GeoPackage
format for subsequent analysis and cartographic
representation. The integration of these
datasets does not, by itself, establish
statistical or causal relationships between
soil properties and vegetation conditions.


\subsection{Sentinel-2 Acquisition and Quality Screening}
\label{subsec:sentinel-acquisition}

Sentinel-2 Level-2A imagery was obtained through
the Copernicus Data Space Ecosystem (CDSE) for
the period November 2025 to June 2026. Products
were discovered through the SpatioTemporal Asset
Catalog (STAC) interface using the geographic
extent and acquisition period relevant to the
study area.

The study area intersects Sentinel-2 tiles
T17QLF and T17QMF. The red band (B04) and
near-infrared band (B08), both provided at
10 m spatial resolution, were selected for
NDVI calculation.

Acquisition records were organized by date,
satellite, and tile to identify complete B04/B08
pairs. Incomplete band pairs were excluded
from NDVI processing.

The Scene Classification Layer (SCL), provided
at 20 m spatial resolution, was used to support
quality screening for potentially cloud- and
shadow-contaminated observations.

The initial satellite inventory contained
226 observations. Following tile-level processing
and preparation of the agricultural analysis
region, 32 acquisition dates were represented
by analysis-area NDVI rasters. A subsequent
regional cloud-screening procedure retained
15 acquisitions meeting the final criterion
of less than 10\% cloud coverage within the
agricultural analysis region.

The cloud-coverage threshold was used as an
acquisition-selection criterion. It does not
guarantee that all pixels in the retained
observations are free from clouds, haze, or
shadows. Cloud screening and numerical NDVI
validation were therefore treated as distinct
quality-control procedures.

Regional cloud coverage was calculated
using SCL classes 8, 9, and 10,
representing medium-probability clouds,
high-probability clouds, and thin cirrus,
respectively.

Cloud shadows, represented by SCL
class 3, were evaluated separately
and were not included in the reported
cloud-coverage percentage.

\subsection{Radiometric Correction and NDVI Calculation}
\label{subsec:ndvi-method}

Sentinel-2 Level-2A spectral bands are distributed
as integer digital numbers (DN). Bottom-of-atmosphere
(BOA) reflectance was calculated using the
radiometric quantification value and additive
offset associated with the satellite products:

\begin{equation}
	\rho_{\mathrm{BOA}} =
	\frac{DN + O_{\mathrm{BOA}}}{Q_{\mathrm{BOA}}},
	\label{eq:boa-reflectance}
\end{equation}

where $DN$ represents the stored pixel value,
$O_{\mathrm{BOA}}$ is the additive radiometric
offset, and $Q_{\mathrm{BOA}}$ is the reflectance
quantification value.

An examination of 111 available Sentinel-2
product XML metadata files identified a
quantification value of 10,000 and additive
offsets of $-1000$ for both B04 and B08.
The examined products corresponded to
processing baselines 05.11 and 05.12.

For consistency, these parameters were applied
to all images used in the analysis, including
products for which XML metadata were unavailable.
Consequently, the radiometric parameters were
not independently verified for every processed
image.

Under this assumption, reflectance was calculated
as:

\begin{equation}
	\rho_{\mathrm{BOA}} =
	\frac{DN - 1000}{10000}.
	\label{eq:boa-correction}
\end{equation}

The Normalized Difference Vegetation Index
(NDVI) was calculated from the corrected red
and near-infrared reflectance values:

\begin{equation}
	\mathrm{NDVI} =
	\frac{\rho_{\mathrm{B08}}-\rho_{\mathrm{B04}}}
	{\rho_{\mathrm{B08}}+\rho_{\mathrm{B04}}},
	\label{eq:ndvi}
\end{equation}

where $\rho_{\mathrm{B08}}$ represents
near-infrared BOA reflectance and
$\rho_{\mathrm{B04}}$ represents red BOA
reflectance.

The additive radiometric offsets were applied
before NDVI calculation. Although a common
multiplicative factor cancels in the normalized
difference, an additive offset does not.
Consequently, the treatment of the additive
offset affects the calculated NDVI values.

Pixels containing zero digital numbers,
nonfinite values, negative corrected reflectance,
or nonpositive reflectance sums were excluded.
Calculated NDVI values outside the theoretical
interval $[-1,1]$ were also rejected.

Invalid pixels were assigned a NoData value
of $-9999$. The resulting NDVI rasters were
exported as single-band, 32-bit floating-point
GeoTIFF files.


\subsection{NDVI Mosaicking and Spatial Extraction}
\label{subsec:ndvi-processing}

NDVI was initially calculated independently
for each complete Sentinel-2 B04/B08 tile pair.

For acquisition dates with both required tiles,
the corresponding tile-level NDVI rasters from
T17QLF and T17QMF were mosaicked to provide
spatial coverage across the analysis region.

The agricultural analysis region was defined
using five polygons from the integrated
soil--vegetation dataset.

Four polygons were associated with
Ferralítico Rojo Típico
(Typical Red Ferrallitic Soil), and
one was associated with
Ferralítico Amarillento Típico
(Typical Yellowish Ferrallitic Soil).

The selected polygons were combined
into a unified analysis geometry
covering approximately 16.22~km$^2$.
This geometry was used consistently
for regional NDVI extraction across
the selected Sentinel-2 acquisitions.

The geometry of the selected agricultural
analysis region was used to define the
extraction area. The relevant vector geometry
was transformed from EPSG:3795 to EPSG:32617,
the coordinate reference system of the
Sentinel-2 rasters, before spatial extraction.

The mosaics were cropped to the analysis
geometry and exported as GeoTIFF files at
10 m spatial resolution.

The processing inventory contained 64
tile-level NDVI rasters and 32 analysis-area
NDVI rasters, corresponding to 32 acquisition
dates. These counts describe the processed
inventory before the final regional cloud
screening used to select the 15 acquisitions
for temporal analysis.


\subsection{QGIS-Based Selection of Potato Fields}
\label{subsec:qgis-potato-fields}

QGIS was used to identify three potato fields
within the processed agricultural region.
The selected locations were verified using
local knowledge of the study area and were
used for a more focused examination of
vegetation-index variability.

The sampling areas were digitized as three
separate circular polygons and overlaid on
the processed NDVI rasters for visual
inspection. These polygons represent
selected sampling areas rather than surveyed
field boundaries.

The same three sampling polygons were used
to extract NDVI observations for the available
acquisition dates, allowing temporal
comparisons within and between the sampled
areas.

The field-level dataset contained 45
field-date observations, corresponding to
three sampling areas evaluated over 15
Sentinel-2 acquisition dates. For each
sampling area and date, the extracted data
were summarized using the number of valid
pixels, mean NDVI, median NDVI, and standard
deviation.

The resulting dataset supports comparisons
of temporal NDVI patterns among the three
sampled areas. Because the sampling polygons
represent selected locations rather than
a comprehensive survey of potato cultivation,
the results should not be interpreted as
representative of every potato field in
the study region.


\subsection{Raster Validation and Temporal Sampling}
\label{subsec:raster-validation}

The analysis-area NDVI rasters were validated
to assess their spatial and numerical
consistency.

Validation included checks of the coordinate
reference system, spatial resolution, raster
dimensions, data type, number of bands,
NoData representation, and range of valid
NDVI values.

The expected raster characteristics were
EPSG:32617, 10 m spatial resolution, a single
32-bit floating-point band, and a NoData
value of $-9999$. Valid NDVI values were
required to fall within the interval $[-1,1]$.

A 15-date temporal sampling design was used
to examine vegetation-index variability across
three broad periods: November--February,
March--April, and May--June. Candidate
acquisitions were evaluated according to
temporal distribution and the availability
of suitable observations.

Only dates represented by processed and
validated imagery were included in the
subsequent temporal analysis.


\subsection{Temporal and Agricultural-Region Analysis}
\label{subsec:temporal-analysis}

The accepted NDVI observations were analyzed
to characterize temporal variation in
vegetation-index values within the agricultural
analysis region and the three selected
sampling areas.

Regional NDVI summaries were used to examine
changes across acquisition dates. Field-level
summaries were used to compare the mean,
median, and variability of NDVI within each
sampling polygon over the same observation
period.

Tables and plots were generated to examine
temporal trends and differences among
the sampled areas. The analysis was
descriptive and focused on the distribution
and evolution of satellite-derived NDVI
values.

The resulting observations provide information
about vegetation-index variability during
the study period. They are not direct
measurements of crop yield, soil properties,
or agricultural management, and they do not
independently establish the causes of
observed temporal changes.